For a finite rank \(k_a\), the elementary order-statistic identity is
\[
s_a(O_{T,a})\le Q_{k_a}(s_a(O_{1,a}),\ldots,s_a(O_{M,a}))
\iff
\sum_{i=1}^M\mathbf 1\{s_a(O_{i,a})<s_a(O_{T,a})\}<k_a.
\tag{1}
\]
Thus ties are handled by strict prefixes and the event depends on \(s_a\) only through its total preorder on \(\mathcal O_a\). Conversely, if a preorder has ordered nonempty classes
\(C_{a,1}\prec\cdots\prec C_{a,d_a}\), assign score zero when \(d_a=1\), and otherwise assign
\((j-1)K/(d_a-1)\) to class \(j\). This realizes it in \([0,K]\) with at most \(m_a\) levels. Hence the finite maximum in
\(\texttt{def:coverage-frontier}\) is attained by the
\(F_{m_0}F_{m_1}\) preorder pairs. This is exactness over the declared score class, not over all invariant latent laws.

When both ranks are finite,
\(\texttt{lem:coherent-two-finite-rank-frontier-evaluator}\) proves
\[
\Pr\{\text{both rank events under }C\}
=\sum_{\pi\in\Pi_C}\omega_C(\pi)D_{M,k_0,k_1}(\pi).
\tag{2}
\]
It preserves the unchanged \(R_w\), coherent \(P_T^C\), local weights, common selector, common target, and exact memo key. The vertex, edge, and zero-corner face prepasses followed by the least-charge nonface dispatcher give exactly the displayed six-way partition. In particular, the face stratum contributes
\(\sum_{\pi\in\Pi_{\mathrm{face}}}\mathsf F_M(\pi)\) to \(\mathsf B\), contributes no nonconstant term to \(\mathsf W\), and does not trigger the multinomial factorial table. The cited evaluator lemma also proves both time formulas and the memory formula. Thus every two-finite-rank computational assertion follows.

The old strict edge audit remains valid: at \(M=2n\), \(k_0=n\), \(k_1=M\), and \(\pi=(1/2,1/2,0,0)\),
\[
D_{M,n,M}(\pi)=1-2^{-M},
\]
with \(\Theta(\log M)\) edge cost versus \(\Theta(M)\) for the former dispatcher; its displayed armwise levels yield the stated ranks. The new face audit in the same lemma gives, for every \(M\ge2\),
\[
k_0=k_1=M,\qquad
\pi=(1/3,1/3,1/3,0),\qquad
D_{M,M,M}(\pi)=1-2\cdot3^{-M}.
\tag{3}
\]
The legal trivial-group two-pair construction produces the three strict-prefix pairs \(00,01,10\) with equal probability. Equation (3) follows because the all-\(10\) and all-\(01\) samples are the two disjoint failures. The face cost is \(\Theta(\log M)\), while rolling, multinomial summation, and every allowed positive-base coefficient orientation cost \(\Theta(M)\). The choice
\(\varepsilon_0=\varepsilon_1=1/(M+1)\) and
\(\alpha=2/(M+1)\) yields both ranks \(M\).

If exactly one rank \(k_b=M+1\), the boundary event is certain by
\(\texttt{prop:asymmetric-rank-boundary}\). For \(c=1-b\), conditioning on the target orbit and applying (1) makes the strict-prefix count binomial with probability
\[
\theta(O_t)=\sum_{O:s_c(O)<s_c(O_t)}q_{c,O}.
\]
Averaging gives the displayed one-arm formula. The cumulative ordered-class scan and
\(\texttt{lem:exact-binomial-shorter-tail-evaluator}\) give the stated
\(F_{m_c}\)-preorder time and constant auxiliary CDF memory, as proved in
\(\texttt{lem:one-boundary-orbit-frontier-evaluator}\). If both ranks are \(M+1\), both events are certain and the positive-part deficit is zero.

For completeness, the binary nonattainment claim is an exact finite audit. In the admissible four-pair instance, take calibration block probabilities
\[
q=(13/20,1/20,1/4,1/20)
\]
and common-target block probabilities
\[
p=(1/20,1/10,1/4,3/5).
\]
At \(M=3\), \(\alpha=1/2\), and equal arm levels \(1/4\), both ranks are three. Taking \(s_0\) constant and ordering arm \(1\) as
\(1\prec3\prec2\prec4\) gives failure probability
\[
\frac14\left(\frac{13}{20}\right)^3+
\frac1{10}\left(\frac{18}{20}\right)^3+
\frac35\left(\frac{19}{20}\right)^3
=\frac{104957}{160000},
\]
and hence objective \(24957/160000\). For binary high-score sets \(H_0,H_1\), direct inclusion--exclusion gives
\[
p(H_0)\{1-q(H_0)\}^3+
p(H_1)\{1-q(H_1)\}^3-
p(H_0\cap H_1)\{1-q(H_0)\}^3\{1-q(H_1)\}^3.
\tag{4}
\]
Order the sixteen choices of \(H_0\) as
\[
\varnothing,1,2,3,4,12,13,14,23,24,34,123,124,134,234,1234.
\]
For these choices, exact maximization of (4) over all sixteen \(H_1\)'s gives, with common denominator \(320000000\), the numerators
\[
\begin{split}
(164616000,&165302000,192052000,198366000,207906936,165912000,164712000,165787368,\\
&203032000,208582976,208582976,164632000,165793875,164631423,202894331,164616000).
\end{split}
\]
The largest failure probability is therefore \(208582976/320000000\), so the largest binary objective is
\(759109/5000000\). Subtraction gives the positive gap
\[
\frac{24957}{160000}-\frac{759109}{5000000}
=\frac{83189}{20000000}.
\]
Thus binary scores are not universally extremal.

Finally, if target and calibration orbit supports are disjoint in a finite-rank arm, put score \(K\) on target-supported orbits and zero on all calibration-supported orbits in that arm, and put zero in the other arm. The adverse rank event fails surely, so the joint event has probability zero and the frontier objective is \(1-\alpha\). No objective can exceed \(1-\alpha\), proving equality. Setting
\(\varepsilon_0=\varepsilon_1=\alpha/2\) gives
\(k_0=k_1=k\) by
\(\texttt{lem:armwise-finite-rank-staircase}\) and yields the diagonal theorem. Every prepass changes only how the same probability in (2) is evaluated, so all active-arm, boundary, pair-audit, empty-library, ordinary-rank, and weighted-orbit semantics are unchanged.